
\documentclass[final]{beamer}
\usepackage[size=a0,scale=1.2]{beamerposter}
\usetheme{ConfPoster}

\usepackage[utf8]{inputenc}
\usepackage[T1]{fontenc}
\usepackage[brazil]{babel}
\usepackage{graphicx}
\usepackage{multicol}
\usepackage{booktabs}

\title{TÍTULO DO SEU PROJETO}
\author{
Nome do Orientador$^{1}$ \\
Nome do Aluno$^{2}$ \\
Nome do Aluno$^{3}$
}
\institute{
$^{1}$Professor do nível tal, Instituição, Município--Estado \\
$^{2}$Discente da série tal, Instituição, Município--Estado \\
$^{3}$Discente da série tal, Instituição, Município--Estado \\
\texttt{email@exemplo.com}
}

\begin{document}
\begin{frame}[t]

\begin{multicols}{3}

\block{Introdução}{
Falar sobre aquilo que moveu o desenvolvimento do projeto, sua importância
e um resumo do embasamento teórico da pesquisa. Logo após, apresentar os objetivos.
}

\block{Objetivo Geral e Específicos}{
\begin{itemize}
  \item Objetivo geral do projeto.
  \item Objetivo específico 1.
  \item Objetivo específico 2.
\end{itemize}
}

\block{Metodologia}{
Descrever cada etapa do processo de pesquisa, desde a parte teórica até o trabalho
de campo. Podem ser utilizados esquemas, fotos e gráficos.
}

\block{Resultados e Conclusões}{
Descrever os resultados obtidos de forma objetiva e discutir os dados observados.
Indicar as conclusões alcançadas, metas atingidas e produtos desenvolvidos.
}

\block{Agradecimentos}{
Inserir agradecimentos às instituições de apoio e fomento.
}

\block{Referências}{
\begin{enumerate}
  \item Koen YM et al. \textit{Chem Res Toxicol}. 2016;29:1064.
  \item Koen YM et al. \textit{Chem Res Toxicol}. 2016;29:1064.
\end{enumerate}
}

\block{Palavras-chave}{
Palavra 1; Palavra 2; Palavra 3.
}

\end{multicols}

\end{frame}
\end{document}
